% Practice Test 18 — Oklahoma Edition
% Questions: 30 (26 MC + 4 SA)
% Generated by scripts/generate_practice_tests.py

\begin{practiceQuestion}{1-3-q12}{mc}
\begin{questionText}
A proportional table has $k = \frac{3}{4}$. Which pair could be in the table?
\end{questionText}
\choiceA{$(8, 10)$}
\choiceB{$(8, 6)$}
\choiceC{$(12, 8)$}
\choiceD{$(4, 2)$}
\correctAnswer{B}
\explanation{If $k = \frac{3}{4}$, then $y = \frac{3}{4}x$. For $x = 8$: $y = \frac{3}{4} \times 8 = 6$. So $(8, 6)$ works.}
\end{practiceQuestion}

\begin{practiceQuestion}{1-6-q04}{mc}
\begin{questionText}
A recipe for $4$ people uses $\frac{3}{2}$ cups of rice. How much rice is needed for $10$ people?
\end{questionText}
\choiceA{$2\frac{1}{2}$ cups}
\choiceB{$3$ cups}
\choiceC{$3\frac{3}{4}$ cups}
\choiceD{$5$ cups}
\correctAnswer{C}
\explanation{Per person: $\frac{3}{2} \div 4 = \frac{3}{8}$ cup. For $10$: $\frac{3}{8} \times 10 = \frac{30}{8} = 3\frac{3}{4}$ cups.}
\end{practiceQuestion}

\begin{practiceQuestion}{1-7-q12}{mc}
\begin{questionText}
A swimming pool is $25$ m long. A scale model is $50$ cm long. What is the scale?
\end{questionText}
\choiceA{$1 : 25$}
\choiceB{$1 : 50$}
\choiceC{$1 : 500$}
\choiceD{$1 : 5$}
\correctAnswer{B}
\explanation{Convert: $25$ m $= 2{,}500$ cm. Scale: $2{,}500 \div 50 = 50$. The scale is $1 : 50$.}
\end{practiceQuestion}

\begin{practiceQuestion}{2-1-q12}{mc}
\begin{questionText}
$63$ is $90\%$ of what number?
\end{questionText}
\choiceA{$56.7$}
\choiceB{$67$}
\choiceC{$70$}
\choiceD{$72$}
\correctAnswer{C}
\explanation{Divide: $63 \div 0.90 = 70$.}
\end{practiceQuestion}

\begin{practiceQuestion}{2-4-q11}{mc}
\begin{questionText}
A $\$200$ television is discounted $30\%$. Then the store applies $5\%$ sales tax. What is the final price?
\end{questionText}
\choiceA{$\$130$}
\choiceB{$\$133$}
\choiceC{$\$140$}
\choiceD{$\$147$}
\correctAnswer{D}
\explanation{Discount: $200 \times 0.70 = \$140$. Tax: $140 \times 1.05 = \$147$.}
\end{practiceQuestion}

\begin{practiceQuestion}{2-7-q02}{mc}
\begin{questionText}
A student measured a table as $4.5$ feet long. The actual length is $5$ feet. What is the percent error?
\end{questionText}
\choiceA{$0.5\%$}
\choiceB{$5\%$}
\choiceC{$10\%$}
\choiceD{$11.1\%$}
\correctAnswer{C}
\explanation{$\frac{|4.5 - 5|}{5} \times 100 = \frac{0.5}{5} \times 100 = 10\%$.}
\end{practiceQuestion}

\begin{practiceQuestion}{2-8-q13}{mc}
\begin{questionText}
Which is a disadvantage of using a credit card?
\end{questionText}
\choiceA{You can buy things before you have the cash}
\choiceB{It builds your credit history}
\choiceC{You pay interest if you don't pay the full balance on time}
\choiceD{It is accepted at most stores}
\correctAnswer{C}
\explanation{The main disadvantage is paying extra money (interest) if you don't pay the full balance by the due date.}
\end{practiceQuestion}

\begin{practiceQuestion}{3-1-q03}{mc}
\begin{questionText}
Which statement is true?
\end{questionText}
\choiceA{$-5 > -2$}
\choiceB{$-3 > 0$}
\choiceC{$-7 < -4$}
\choiceD{$0 < -1$}
\correctAnswer{C}
\explanation{On a number line, $-7$ is to the left of $-4$, so $-7 < -4$.}
\end{practiceQuestion}

\begin{practiceQuestion}{4-1-q14}{mc}
\begin{questionText}
What is the value of $x^2 + 3x$ when $x = -2$?
\end{questionText}
\choiceA{$10$}
\choiceB{$-2$}
\choiceC{$-10$}
\choiceD{$2$}
\correctAnswer{B}
\explanation{Substitute: $(-2)^2 + 3(-2) = 4 + (-6) = -2$.}
\end{practiceQuestion}

\begin{practiceQuestion}{4-2-q12}{mc}
\begin{questionText}
Simplify $-6y + 4y - y$.
\end{questionText}
\choiceA{$-3y$}
\choiceB{$-y$}
\choiceC{$3y$}
\choiceD{$-11y$}
\correctAnswer{A}
\explanation{$-6y + 4y = -2y$. Then $-2y - y = -3y$.}
\end{practiceQuestion}

\begin{practiceQuestion}{5-3-q11}{mc}
\begin{questionText}
A student solved $3(x + 4) = 30$ and wrote: $3x + 4 = 30$, $3x = 26$, $x = \frac{26}{3}$. What was the student's error?
\end{questionText}
\choiceA{Divided by $3$ instead of subtracting}
\choiceB{Did not distribute $3$ to the $4$}
\choiceC{Subtracted $4$ from the wrong side}
\choiceD{Forgot to check the answer}
\correctAnswer{B}
\explanation{The student wrote $3x + 4$ instead of $3x + 12$. Distributing correctly gives $3x + 12 = 30$.}
\end{practiceQuestion}

\begin{practiceQuestion}{5-4-q25}{sa}
\begin{questionText}
Solve $2.5x - 7.5 = 10$.
\end{questionText}
\correctAnswer{$x = 7$}
\explanation{Add $7.5$: $2.5x = 17.5$. Divide by $2.5$: $x = 7$. Check: $2.5(7) - 7.5 = 17.5 - 7.5 = 10$ \checkmark.}
\end{practiceQuestion}

\begin{practiceQuestion}{5-5-q06}{mc}
\begin{questionText}
Solve $-4x > 20$.
\end{questionText}
\choiceA{$x > -5$}
\choiceB{$x < -5$}
\choiceC{$x > 5$}
\choiceD{$x < 5$}
\correctAnswer{B}
\explanation{Divide by $-4$ and flip the sign: $x < -5$.}
\end{practiceQuestion}

\begin{practiceQuestion}{5-6-q20}{sa}
\begin{questionText}
Solve $-5x + 15 \geq 30$. Describe the graph.
\end{questionText}
\correctAnswer{$x \leq -3$; closed circle at $-3$, shade left}
\explanation{Subtract $15$: $-5x \geq 15$. Divide by $-5$ and flip: $x \leq -3$. Closed circle at $-3$, shade left.}
\end{practiceQuestion}

\begin{practiceQuestion}{6-2-q14}{mc}
\begin{questionText}
Original scale: $1$ cm $= 10$ m. New scale: $1$ cm $= 40$ m. What is the scale ratio?
\end{questionText}
\choiceA{$\frac{1}{4}$}
\choiceB{$4$}
\choiceC{$30$}
\choiceD{$\frac{1}{30}$}
\correctAnswer{A}
\explanation{Scale ratio: $\frac{10}{40} = \frac{1}{4}$. Every old length is multiplied by $\frac{1}{4}$ on the new drawing.}
\end{practiceQuestion}

\begin{practiceQuestion}{6-3-q17}{mc}
\begin{questionText}
A triangle has side lengths $9$, $12$, and $15$. What type of triangle is it?
\end{questionText}
\choiceA{Equilateral}
\choiceB{Isosceles}
\choiceC{Right}
\choiceD{It cannot be formed}
\correctAnswer{C}
\explanation{Check: $9^2 + 12^2 = 81 + 144 = 225 = 15^2$. Since the Pythagorean theorem holds, this is a right triangle.}
\end{practiceQuestion}

\begin{practiceQuestion}{6-4-q27}{gmc}
\begin{questionText}
The table below shows four sets of triangle conditions. Which one produces a unique triangle?

\begin{center}
\begin{tabular}{|c|l|l|}
\hline
\textbf{Set} & \textbf{Given Information} & \textbf{Condition} \\
\hline
A & Angles: $50°$, $60°$, $70°$ & AAA \\
\hline
B & Sides: $3$, $4$, $5$ & SSS \\
\hline
C & Angles: $45°$, $90°$, $45°$ & AAA \\
\hline
D & One side: $8$ cm & --- \\
\hline
\end{tabular}
\end{center}
\end{questionText}
\choiceA{Set A}
\choiceB{Set B}
\choiceC{Set C}
\choiceD{Set D}
\correctAnswer{B}
\explanation{SSS with valid side lengths ($3 + 4 = 7 > 5$, etc.) produces exactly one unique triangle. AAA gives many triangles, and one side alone is insufficient.}
\end{practiceQuestion}

\begin{practiceQuestion}{6-5-q09}{mc}
\begin{questionText}
You cut a cone with a vertical slice through its apex (tip). What shape is the cross-section?
\end{questionText}
\choiceA{Circle}
\choiceB{Triangle}
\choiceC{Rectangle}
\choiceD{Semicircle}
\correctAnswer{B}
\explanation{A vertical cut through the apex of a cone produces a triangle (specifically an isosceles triangle).}
\end{practiceQuestion}

\begin{practiceQuestion}{6-6-q17}{mc}
\begin{questionText}
An angle is $12°$ less than its complement. What is the angle?
\end{questionText}
\choiceA{$39°$}
\choiceB{$42°$}
\choiceC{$48°$}
\choiceD{$51°$}
\correctAnswer{A}
\explanation{Let the angle be $x$ and its complement be $x + 12$. Sum: $x + (x + 12) = 90$. Combine: $2x + 12 = 90$. Solve: $2x = 78$, $x = 39°$.}
\end{practiceQuestion}

\begin{practiceQuestion}{7-4-q11}{mc}
\begin{questionText}
A T-shaped figure can be broken into a $12$ cm by $2$ cm horizontal rectangle and a $2$ cm by $8$ cm vertical rectangle. What is the total area?
\end{questionText}
\choiceA{$24$ cm$^2$}
\choiceB{$32$ cm$^2$}
\choiceC{$40$ cm$^2$}
\choiceD{$48$ cm$^2$}
\correctAnswer{C}
\explanation{Horizontal: $12 \times 2 = 24$ cm$^2$. Vertical: $2 \times 8 = 16$ cm$^2$. Total: $24 + 16 = 40$ cm$^2$.}
\end{practiceQuestion}

\begin{practiceQuestion}{7-5-q16}{mc}
\begin{questionText}
If all dimensions of a rectangular prism are doubled, what happens to its surface area?
\end{questionText}
\choiceA{It doubles}
\choiceB{It triples}
\choiceC{It quadruples}
\choiceD{It increases by a factor of $8$}
\correctAnswer{C}
\explanation{Surface area depends on the square of the dimensions. Doubling all dimensions multiplies each face area by $2^2 = 4$, so the total surface area quadruples.}
\end{practiceQuestion}

\begin{practiceQuestion}{7-6-q29}{gsa}
\begin{questionText}
A composite solid is made of two rectangular prisms joined together as shown. What is the total volume?

\begin{center}
\begin{tikzpicture}[scale=0.35]
  % Bottom prism
  \draw[thick,funBlue,fill=funBlue!10] (0,0) rectangle (10,3);
  \draw[thick,funBlue,fill=funBlue!20] (10,0) -- (13,2) -- (13,5) -- (10,3) -- cycle;
  \draw[thick,funBlue,fill=funBlue!15] (0,3) -- (3,5) -- (13,5) -- (10,3) -- cycle;
  % Top prism
  \draw[thick,funRed,fill=funRed!10] (0,3) rectangle (4,6);
  \draw[thick,funRed,fill=funRed!20] (4,3) -- (7,5) -- (7,8) -- (4,6) -- cycle;
  \draw[thick,funRed,fill=funRed!15] (0,6) -- (3,8) -- (7,8) -- (4,6) -- cycle;
  % Labels
  \node[below] at (5,0) {$10$ cm};
  \node[left] at (0,1.5) {$3$ cm};
  \node[right] at (13,3.5) {$3$ cm};
  \node[left] at (0,4.5) {$3$ cm};
  \node[above left] at (2,6) {$4$ cm};
  \node[right] at (7,6.5) {$3$ cm};
\end{tikzpicture}
\end{center}
\end{questionText}
\correctAnswer{$126$ cm$^3$}
\explanation{Bottom prism: $10 \times 3 \times 3 = 90$ cm$^3$. Top prism: $4 \times 3 \times 3 = 36$ cm$^3$. Total: $90 + 36 = 126$ cm$^3$.}
\end{practiceQuestion}

\begin{practiceQuestion}{8-1-q18}{mc}
\begin{questionText}
A poll found that $35\%$ of a random sample of $400$ students prefer pizza for lunch. Which statement is the best inference?
\end{questionText}
\choiceA{Exactly $35\%$ of all students prefer pizza}
\choiceB{About $35\%$ of the entire student population likely prefers pizza}
\choiceC{No students prefer anything other than pizza}
\choiceD{The poll is invalid because not all students were surveyed}
\correctAnswer{B}
\explanation{A random sample provides an estimate. About $35\%$ of the population likely prefers pizza, but the exact percentage may differ slightly.}
\end{practiceQuestion}

\begin{practiceQuestion}{8-3-q12}{mc}
\begin{questionText}
When is a dot plot a better choice than a box plot for comparing two groups?
\end{questionText}
\choiceA{When data sets are very large}
\choiceB{When you want to see every individual data point}
\choiceC{When you only want to compare medians}
\choiceD{When the data is categorical}
\correctAnswer{B}
\explanation{Dot plots show every data point, which is useful for small data sets where individual values matter.}
\end{practiceQuestion}

\begin{practiceQuestion}{8-4-q18}{mc}
\begin{questionText}
Class A: median $85$, IQR $8$, range $20$. Class B: median $80$, IQR $15$, range $40$. A student says ``Class B is better because it has a wider range.'' Is the student correct?
\end{questionText}
\choiceA{Yes — wider range means more talent}
\choiceB{No — wider range means less consistency, and Class A has a higher median}
\choiceC{Yes — Class B's scores spread more, which is always better}
\choiceD{No — range does not measure anything useful}
\correctAnswer{B}
\explanation{A wider range indicates more variability, not better performance. Class A has a higher median and is more consistent (smaller IQR and range).}
\end{practiceQuestion}

\begin{practiceQuestion}{8-5-q22}{sa}
\begin{questionText}
How many data values are shown in this stem-and-leaf plot?

Stem $5$: $1\;\;3\;\;5$ \quad Stem $6$: $0\;\;2\;\;4\;\;7\;\;8$ \quad Stem $7$: $1\;\;6$ \quad Stem $8$: $0$
\end{questionText}
\correctAnswer{$11$}
\explanation{Count leaves: $3 + 5 + 2 + 1 = 11$.}
\end{practiceQuestion}

\begin{practiceQuestion}{8-7-q18}{mc}
\begin{questionText}
Which statement is true about frequency tables?
\end{questionText}
\choiceA{They only work for numerical data}
\choiceB{They can organize both categorical and numerical data}
\choiceC{They always require intervals}
\choiceD{They show data as a graph}
\correctAnswer{B}
\explanation{Frequency tables can list categories (favorite color) or numerical intervals (test scores $70$--$79$).}
\end{practiceQuestion}

\begin{practiceQuestion}{9-4-q03}{mc}
\begin{questionText}
A bag contains $4$ red marbles and $1$ blue marble. Which type of probability model does this represent?
\end{questionText}
\choiceA{Uniform}
\choiceB{Non-uniform}
\choiceC{Both uniform and non-uniform}
\choiceD{Neither}
\correctAnswer{B}
\explanation{The probability of drawing red ($\frac{4}{5}$) is different from drawing blue ($\frac{1}{5}$), so this is a non-uniform model.}
\end{practiceQuestion}

\begin{practiceQuestion}{9-6-q04}{mc}
\begin{questionText}
Two dice are rolled. What is the probability of getting a sum of $2$?
\end{questionText}
\choiceA{$\frac{1}{6}$}
\choiceB{$\frac{1}{12}$}
\choiceC{$\frac{1}{36}$}
\choiceD{$\frac{2}{36}$}
\correctAnswer{C}
\explanation{The only way to get a sum of $2$ is $(1,1)$. $P = \frac{1}{36}$.}
\end{practiceQuestion}

\begin{practiceQuestion}{9-7-q15}{mc}
\begin{questionText}
You simulate flipping a coin $3$ times and counting how many times all $3$ are heads. In $40$ simulations, all $3$ heads occurs $6$ times. What is the experimental probability?
\end{questionText}
\choiceA{$\frac{6}{40} = 0.15$}
\choiceB{$\frac{1}{8} = 0.125$}
\choiceC{$\frac{3}{40} = 0.075$}
\choiceD{$\frac{6}{120} = 0.05$}
\correctAnswer{A}
\explanation{$P = \frac{6}{40} = 0.15$. The theoretical probability is $\frac{1}{8} = 0.125$, close to the simulation result.}
\end{practiceQuestion}
